\begin{lemma}
\label{lemma-minimal-contains}
Let $M$ be a flat Mittag-Leffler module over $R$. Let $F$ be an $R$-module
and let $x \in F \otimes_R M$. Then there exists a smallest submodule
$F' \subset F$ such that $x \in F' \otimes_R M$.
Also, $F'$ is a finite $R$-module.
\end{lemma}

\begin{proof}
Since $M$ is flat we have $F' \otimes_R M \subset F \otimes_R M$
if $F' \subset F$ is a submodule, hence the statement makes sense.
Let $I = \{F' \subset F \mid x \in F' \otimes_R M\}$ and for
$i \in I$ denote $F_i \subset F$ the corresponding submodule.
Then $x$ maps to zero under the map
$$
F \otimes_R M \longrightarrow \prod (F/F_i \otimes_R M)
$$
whence by Proposition \ref{proposition-ML-tensor}
$x$ maps to zero under the map
$$
F \otimes_R M \longrightarrow \left(\prod F/F_i\right) \otimes_R M
$$
Since $M$ is flat the kernel of this arrow is
$(\bigcap F_i) \otimes_R M$ which proves that $F' = \bigcap F_i$.
To see that $F'$ is a finite module, suppose that
$x = \sum_{j = 1, \ldots, m} f_j \otimes m_j$ with $f_j \in F'$
and $m_j \in M$. Then $x \in F'' \otimes_R M$ where $F'' \subset F'$ is
the submodule generated by $f_1, \ldots, f_m$. Of course then
$F'' = F'$ and we conclude the final statement holds.
\end{proof}
